\chapter{Roteiro para Projeto de TCC}

Este roteiro pode ser utilizado como apoio para elaboração do projeto de TCC. O curso pode adaptar, reduzir ou ampliar os itens conforme sua área.

O roteiro segue a lógica da ABNT NBR 15287:2025, mas não substitui o formulário ou as exigências aprovadas pelo curso.

\section{Elementos pré-textuais}

Utilize folha de rosto com autoria, título, natureza do projeto, vínculo acadêmico, orientação, local e ano. Inclua capa, listas e sumário quando forem exigidos ou necessários. Não antecipe no projeto elementos próprios da versão final, como folha de aprovação.

\section{Identificação}

\begin{itemize}
  \item Nome do estudante.
  \item Curso.
  \item Orientador.
  \item Título provisório.
  \item Linha de pesquisa ou área, quando houver.
\end{itemize}

\section{Tema e delimitação}

Descreva o tema geral e o recorte adotado. Indique espaço, tempo, população, material, objeto ou enfoque. A delimitação deve mostrar o que será estudado e o que ficará fora do trabalho.

\section{Problema de pesquisa}

Apresente a pergunta central do trabalho. Ela deve ser clara, investigável e compatível com o método.

\section{Objetivos}

Apresente um objetivo geral e objetivos específicos. Os objetivos específicos devem indicar etapas intelectuais necessárias para alcançar o objetivo geral.

\section{Justificativa}

Explique a relevância do trabalho. Considere importância acadêmica, social, profissional, institucional ou tecnológica.

\section{Referencial teórico inicial}

Apresente autores, conceitos e debates que sustentam a pesquisa. Mostre como essas fontes ajudam a compreender o problema.

\section{Metodologia}

Indique tipo de pesquisa, abordagem, materiais, participantes, instrumentos, procedimentos, critérios de análise e cuidados éticos.

\section{Recursos e riscos}

Registre recursos necessários, formas de acesso e riscos capazes de afetar a execução. Indique medidas alternativas para autorizações, recrutamento, coleta, equipamentos ou bases essenciais.

\section{Cronograma}

\begin{longtable}{p{0.45\linewidth}p{0.37\linewidth}}
\toprule
Atividade & Período previsto \\
\midrule
\endfirsthead
\toprule
Atividade & Período previsto \\
\midrule
\endhead
Delimitação do tema e problema & \\
Levantamento bibliográfico & \\
Elaboração do projeto & \\
Submissão ética, quando aplicável & \\
Coleta ou seleção de dados & \\
Análise & \\
Redação das seções textuais & \\
Revisão com orientador & \\
Normalização & \\
Entrega para banca & \\
Correções finais & \\
\bottomrule
\end{longtable}

\section{Referências iniciais}

Liste as fontes que sustentam a proposta. O projeto não precisa conter bibliografia definitiva, mas deve demonstrar base mínima para início da pesquisa.

\section{Apêndices, anexos e índice}

Inclua instrumentos elaborados pela autoria como apêndices e documentos de terceiros como anexos, somente quando necessários. O índice é opcional e deve seguir a ABNT NBR 6034:2004 quando utilizado.
